\scanpage{5}

\mspage{3}

c) Les applications

\[ \Phi\unsure{\mapstochar\relbar|\mapsto}\Phi_*,\qquad\Psi\mapstochar\relbar\Psi^* \]

établissent des bijections \unsure{inv.|inv} l'une de l'autre entre l'ens \deleted{\unsure{$\Phi$}}des parties $\Phi$ closes à gauche de $\mathrm{Fl}\,M$, et l'ens des parties $\Psi$ closes à droite. \inserted{C'est un anti-isom. d'ens. ordonnés.} On a alors

\[ \begin{aligned}(\Phi\cap\Phi')_{\unsure{*|\alpha}}={}&\text{partie close à droite de }\mathrm{Fl}(M)\\ &\text{``engendrée par }\Phi_*,\Phi'_* \text{'', i.e plus}\\ &\text{petite partie close à droite}\\ &\text{contenant }\Phi_*,\Phi'_*,\ \text{i.e.}\\ &\overline{\Phi_*\cup\Phi'_*}\end{aligned} \]

plus généralement

\[ \left(\bigcap_i\Phi_i\right)_*=\deleted{\text{\unsure{partie}}}\overline{\bigcup_i(\Phi_{i*})} \]

pour tte famille $(\Phi_i)_{i\in\unsure{I|J}}$ de parties closes à gauche de $\mathrm{Fl}\,M$. Formules duales

\[ \left\{\begin{aligned}(\Psi\cap\Psi')^*&=\widetilde{\Psi^*\cup\Psi'{}^*},\\ \left(\bigcap_j\Psi_j\right)^*&=\widetilde{\bigcup_j(\Psi_{j\unsure{*}}^{*})}.\end{aligned}\right. \]

d) Si $\Phi,\Psi\subset\mathrm{Fl}(M)$, alors on a

\[ \Phi\subset\Psi^*\Longleftrightarrow\Psi\subset\Phi_*\qquad\text{on écrit alors }\Phi\longleftrightarrow\Psi, \]

cela signifie que $\Phi$ a la LLP pour $\Psi$, ou que $\Psi$ a \unsure{le|la} RLP pour $\Phi$.

% archive sheet number: 4

